\clearpage
\section[Technical Description and Implementation]{Chapter -- 02: Technical Description and Implementation}

This chapter presents the end-to-end pipeline used to load the ASCAD traces, build
per-byte labels, train four families of model, and score each one with the Key
Rank / Guessing Entropy metric.

\subsection{Actual Work}

The pipeline reads the ASCAD HDF5 files into profiling and attack trace sets,
standardises the traces using statistics fitted on the profiling set only, and builds
labels for the chosen key byte. A single command-line entry point,
\texttt{src/run\_experiment.py}, drives every model family through one shared
interface (\texttt{fit} / \texttt{predict\_proba} / \texttt{save} / \texttt{load}), so
that a model can be swapped with a single flag:

\begin{verbatim}
python -m src.run_experiment --model random_forest
python -m src.run_experiment --model svm
python -m src.run_experiment --model cnn
python -m src.run_experiment --model resnet
python -m src.run_experiment --model multitask_resnet
\end{verbatim}

Each run writes a resolved configuration, a full log, the fitted model, a
training-history plot, the key-rank curve (CSV and PNG), and a metrics file to its
own timestamped directory under \texttt{results/}.

\subsection{Prerequisites}

\noindent\textbf{Software:} Python 3.14, PyTorch 2.13.0+cu126 (CUDA 12.6), scikit-learn,
NumPy/SciPy, h5py, Matplotlib.

\noindent\textbf{Hardware:} Windows 11 workstation, NVIDIA RTX 4060 (8\,GB VRAM).
CUDA is optional -- the CNN and ResNet also train on CPU, more slowly.

\noindent\textbf{Technical knowledge:} PyTorch training loops and checkpointing,
convolutional and residual network design, statistical evaluation (guessing entropy),
and HDF5 data handling.

\subsection{Target and Leakage Model}

The attacked value is $y=\mathrm{Sbox}(p[b]\oplus k[b])$ for a chosen key byte $b$,
where $p$ is the known plaintext. Under the identity leakage model this gives 256
classes. Labels are re-derived directly from the trace metadata and cross-checked
against ASCAD's own stored labels for byte 2, which they must match exactly:

\begin{verbatim}
labels[i] = Sbox[ plaintext[i, byte] ^ key[i, byte] ]   # identity model, 256 classes
\end{verbatim}

The attack set is given no single label: every one of the 256 key hypotheses implies
a different label for the same trace, and enumerating all 256 is exactly what the Key
Rank step below does.

\subsection{Model-Selection Fix (Engineering Contribution)}

Validation cross-entropy loss is a poor stopping criterion for masked side-channel
analysis: for roughly the first 20--25 epochs it is flat or even rising while the
attack is already improving, so the lowest-loss epoch corresponds to a
near-untrained network. Every deep model in this project is therefore checkpointed on
\emph{validation guessing entropy}, computed every few epochs on a held-out slice of
the profiling set, rather than on validation loss. This single change is what makes
the ResNet and multi-task ResNet results in Chapter~3 reproducible; selecting on loss
alone silently discards the epoch at which the model has actually learned to attack.
Training also keeps the full trace tensor resident on the GPU, which is several times
faster than per-batch host-to-device transfer for these small models on large trace
sets.

\subsection{Methods Compared}

\begin{itemize}
\item \textbf{Random forest} -- 300 trees, the classical baseline.
\item \textbf{SVM} -- RBF kernel with PCA to 50 components; profiling subsampled to
      10\,000 traces because an RBF SVM is quadratic in sample count.
\item \textbf{CNN} -- the compact ASCAD network of Zaid et al.~\cite{zaid2020} ($\sim$17k
      parameters, SELU activations, one-cycle learning rate). The larger VGG-style
      \texttt{ASCAD\_best} network was also tried but memorises the profiling set
      without ever ranking the key.
\item \textbf{ResNet} -- a 1-D residual network in the style of Karayalcin et
      al.~\cite{karayalcin2023}: kernel size 11, SELU with LeCun-normal
      initialisation, a number of residual blocks scaled as
      $\lfloor\log_2(\text{length})\rfloor-2$, a flatten head, and a triangular
      cyclic learning rate.
\item \textbf{Multi-task ResNet} -- following Marquet and
      Oswald~\cite{marquet2023}, the same trunk with three joint linear heads
      trained simultaneously: $y=\mathrm{Sbox}(p\oplus k)$, the output mask share
      $r_{\text{out}}$, and $y\oplus r_{\text{out}}$. Only the $y$ head is used for
      key recovery; the other two give back-propagation a non-flat gradient during
      the early training plateau.
\end{itemize}

\subsection{The Key Rank / Guessing Entropy Metric}

For $N$ attack traces and every key hypothesis $g\in\{0,\dots,255\}$, the
log-probabilities the model assigns to $\mathrm{Sbox}(p_i\oplus g)$ are summed:

\begin{equation}
\mathrm{score}(g) = \sum_{i=1}^{N} \log P_{\text{model}}\bigl(\text{class} =
\mathrm{Sbox}(p_i \oplus g) \;\big|\; \text{trace}_i\bigr)
\label{eq:keyrank}
\end{equation}

The 256 hypotheses are ranked by this score and the rank of the true key byte is
recorded; rank 0 means the key byte is recovered. This is averaged over 100 random
attack-trace orderings, and the curve of mean rank against number of attack traces is
the primary result reported in Chapter~3.

\subsection{Responsibilities}

\begin{itemize}
\item Building the ASCAD data loader, amplitude scaler, and label pipeline from the
      raw HDF5 files.
\item Implementing the shared model interface and all five model families (random
      forest, SVM, CNN, ResNet, multi-task ResNet).
\item Implementing the Key Rank / Guessing Entropy metric and the validation-GE
      checkpointing fix.
\item Running the fixed-key, variable-key, desynchronised, and shuffled-label
      experiment suites and producing every result figure and table.
\item Implementing the multi-byte assembly step (\texttt{src/assemble\_key.py}) and
      documenting its data limitation.
\end{itemize}

\subsection{Challenges}

\begin{itemize}
\item An early kernel-3, ReLU ResNet overfit the fixed-key set to chance and would
      not train on the variable-key set at all; it took larger kernels, SELU,
      depth scaled to the input length, and a cyclic learning rate to fix (see
      \texttt{docs/resnet\_improvement\_research.md}).
\item Validation cross-entropy is anti-correlated with attack success during the
      first $\sim$20 epochs, which silently breaks naive early stopping (addressed
      above).
\item Deep-model results are seed-sensitive: the variable-key ResNet needed 276
      traces on one seed and 1540 on another, and the desync-50 ResNet needed 80
      epochs rather than 50 to break its training plateau on a repeated run.
\item Full 16-byte key recovery is blocked by the published ASCAD files themselves,
      which are a trace window around byte 2's leakage only; bytes 3 and 4 never
      leave chance when attacked directly (Section~3.6).
\end{itemize}

\subsection{Software Tools Used}

Table~\ref{tab:tools} lists every tool used, so the pipeline can be reproduced
exactly.

\begin{table}[h]
\centering
\caption{Software tools used in the implementation.}
\label{tab:tools}
\begin{tabular}{@{}lll@{}}
\toprule
Software & Version & Purpose \\
\midrule
Python & 3.14 & Core implementation language \\
PyTorch & 2.13.0+cu126 & Model training, GPU acceleration \\
scikit-learn & latest & Random forest, PCA, SVM \\
NumPy / SciPy & latest & Numerical computation \\
h5py & latest & Reading the ASCAD HDF5 files \\
Matplotlib & latest & All result figures (Chapter~3) \\
pytest & latest & Synthetic-leakage pipeline tests \\
\bottomrule
\end{tabular}
\end{table}

\subsection{Pipeline Algorithm}

\begin{algorithm}[h]
\caption{ASCAD Key-Recovery Benchmark Pipeline}
\begin{algorithmic}
\State Load ASCAD HDF5 file; split into profiling and attack trace sets
\State Fit amplitude scaler on profiling traces only; apply to both sets
\State Build labels $y_i = \mathrm{Sbox}(p_i[b] \oplus k_i[b])$ for target byte $b$; cross-check against stored labels
\For{each model $m \in \{$RF, SVM, CNN, ResNet, Multi-task ResNet$\}$}
    \State Train $m$ on profiling traces, checkpointing on validation guessing entropy
    \For{each attack-trace ordering (100 repeats)}
        \State Accumulate $\mathrm{score}(g)$ for all 256 key hypotheses via Eq.~\eqref{eq:keyrank}
        \State Record rank of the true key byte
    \EndFor
    \State Average rank curves; write \texttt{key\_rank.csv} / \texttt{.png} / metrics
\EndFor
\State Emit comparison figures and the summary table for Chapter~3
\end{algorithmic}
\end{algorithm}

\textbf{Result:} The full pipeline -- data loading, five model families, the Key Rank
metric, and the validation-GE checkpointing fix -- was implemented and executed
end-to-end, producing the result tables and figures presented in Chapter~3.
